\documentclass[letterpaper,11pt]{moderncv}
\moderncvstyle{banking}
\usepackage{xpatch}
\moderncvcolor{black} % Subtle, classic color
\AtBeginDocument{\definecolor{color2}{rgb}{0,0,0}} % Force first name black
\usepackage[letterpaper, margin=0.6in]{geometry}
\usepackage{xcolor} % Ensure color support
\usepackage{graphicx} % Support for images
\usepackage{fontspec} % Support for Unicode characters
\setmainfont{Times New Roman} % Set English font
\pagenumbering{arabic} % Ensure page numbering is defined
\usepackage[unicode]{hyperref} % Load hyperref after setting page numbering
\sloppy % Allow LaTeX to relax spacing issues
\usepackage{enumitem} % Allows list customization

\usepackage[english]{babel}
\usepackage[autostyle, english = american]{csquotes}
\MakeOuterQuote{"}

%==================== PERSONAL INFO ====================%
\name{\textcolor{black}{Saeyoung}}{Kim}
\address{Cambridge, MA, USA}
\phone[mobile]{+1 (351) 322 5510}
\email{saeyoung\_kim@uml.edu}
\social[linkedin]{saeyoung-macx-kim}
\social[github]{PrimaryAnomaly}

\begin{document}

\makecvtitle
\vspace{-3em}

%==================== SUMMARY ====================%
\section{Professional Summary}
Applications-oriented engineer with a PhD in Electrical Engineering and hands-on experience spanning MEMS microfabrication, multi-modal sensor development, embedded systems, and statistical process control. Built and validated complex measurement systems from concept through integration testing across space hardware and pharmaceutical manufacturing programs. Skilled in data-driven root cause analysis, scripting-based diagnostics (Python, MATLAB, C/C++), cross-functional stakeholder coordination, and translating customer requirements into actionable technical solutions. Background in cleanroom fabrication processes including photolithography, thin-film deposition, and electroplating.

\section{Technical Skills}
\begin{minipage}[t]{0.48\textwidth}
\cvitem{Process}{Photolithography, Masking, Electroplating, Thin-Film Deposition, MEMS Fabrication}
\cvitem{Analysis}{Statistical Process Control (MSPC), Data Pipeline Design, Root Cause Analysis}
\cvitem{Programming}{Python, C/C++, MATLAB, Scripting \& Automation}
\end{minipage}%
\hfill
\begin{minipage}[t]{0.48\textwidth}
\cvitem{Metrology}{Electrochemical Sensors, Optical Sensing, Multi-Modal Measurement Systems}
\cvitem{Integration}{System-Level V\&V, AIT, Test Plan Development, EGSE Architecture}
\cvitem{Tools}{Git, Docker, Linux, KiCad, Fusion360, AutoCAD}
\end{minipage}

%==================== EXPERIENCE ====================%
\section{Professional Experience}

\cventry{2025.07--Present}{\textbf{Postdoctoral Associate} | Process Modeling and Statistical Process Control}{University of Massachusetts Lowell}{Lowell, MA, USA}{}{
\begin{itemize}[leftmargin=1cm, itemsep=0pt, parsep=0pt]
    \item Developed digital twin framework integrating physics-based modeling with multivariate statistical process control (MSPC) for real-time pharmaceutical manufacturing monitoring
    \item Analyzed 100+ production batches to identify process sensitivity and root causes of quality excursions
    \item Built Python-based analytics suite replacing commercial platforms, enabling flexible scripting and rapid iteration on diagnostic workflows
    \item Designed end-to-end data pipelines for multi-day batch cycles with automated anomaly detection and quality prediction
    \item Translated cross-functional stakeholder requirements into reproducible technical deliverables for government-funded collaboration
    \item Co-authoring journal manuscripts (IF 10+) on process optimization and digital twin methodologies
\end{itemize}}

\cventry{2024.02--2025.05}{\textbf{Systems Integration Engineer} | Satellite Integration, Testing, and Customer Delivery}{Hanwha Systems}{Seoul, South Korea}{}{
\begin{itemize}[leftmargin=1cm, itemsep=0pt, parsep=0pt]
    \item Led system integration and validation of military SAR satellites across EM, QM, and FM test campaigns
    \item Generated and executed comprehensive Integrated Systems Test (IST) plans for Flight Model acceptance testing
    \item Designed Electrical Ground Support Equipment (EGSE) architecture and coordinated integration test configurations across subsystems
    \item Diagnosed and resolved complex system-level anomalies during integration, driving service order closure from a technical perspective
    \item Defined subsystem-level requirements and translated program-level specifications into testable acceptance criteria
    \item Coordinated across subsystem engineers, program office, and external stakeholders to align priorities and resolve escalations
    \item Created knowledge-sharing documentation and test procedures for cross-team reference and onboarding
    \item Contributed to architecture definition for a high-throughput satellite manufacturing facility
\end{itemize}}

%==================== EDUCATION ====================%
\section{Education}
\cventry{2017.03--2024.02}{\textbf{PhD in Electrical Engineering} | MEMS Sensor Systems \& Microfabrication}{Korea University}{Seoul, South Korea}{}{
\begin{itemize}[leftmargin=1cm, itemsep=0pt, parsep=0pt]
    \item Secured and managed \$220K research grant for CubeSat biological payload sensor system development
    \item Designed, fabricated, and validated multi-modal electrochemical and optical sensor suite through full cleanroom process: photolithography, masking, gold electroplating, thin-film deposition, and packaging
    \item Developed embedded firmware in C/C++ for real-time multi-channel data acquisition and autonomous operation
    \item Executed system-level verification and validation of pressurized fluidic and thermal management subsystems
    \item Identified process sensitivities and optimized fabrication parameters through iterative characterization and data analysis
    \item Published 3 first-author SCI papers (cumulative IF: 17.4) and co-authored 10+ publications including work in \textit{Semiconductor Science and Technology}
\end{itemize}}

\cventry{2013.03--2017.02}{\textbf{BSc in Mechanical Engineering}}{Konkuk University}{Seoul, South Korea}{}{
\begin{itemize}[leftmargin=1cm, itemsep=0pt, parsep=0pt]
    \item Conducted research in electrochemical sensor microfabrication, cleanroom processing, and hardware prototyping
    \item Gained experience in CAD modeling, PCB design, precision metrology, and technical documentation
\end{itemize}}

%==================== PUBLICATIONS ====================%
\section{Selected Publications}
\begin{sloppypar}
\begin{itemize}[leftmargin=1cm, itemsep=0pt, parsep=0pt]
    \item \textbf{\underline{Saeyoung Kim}}, Jaewon Lee, Kimia Babaei, Seonkgyu Yoon, "Digital Twins in Lyophilization: Advances in Pharmaceutical Process Development and Optimization," \textit{Biotechnology Advances}, 2026. (In Preparation)
    \item \textbf{\underline{Kim, S.}}; Park, S.; Pak, J.J., "Multi-Modal Multi-Array Electrochemical and Optical Sensor Suite for a Biological CubeSat Payload," \textit{MDPI Sensors}, 2024, 24, Art. no. 265.
    \item \textbf{\underline{Saeyoung Kim}}, Ji-Hoon Han, Beelee Chua, James Jungho Pak, "A pH Sensing Pipette for Cross-Contamination Prevention in Industrial Fermentation," \textit{IEEE Trans. Industrial Electronics}, 2022, 69(7), 7461-7469.
    \item Sanghoon Cho, Jungmo Jung, \textbf{\underline{Saeyoung Kim}}, James Pak, "Conduction Mechanism and Synaptic Behaviour of Interfacial Switching AlO$\sigma$-based RRAM," \textit{Semiconductor Science and Technology}, 2020, 35(8), 085006.
    \item Ji-Hoon Han, \textbf{\underline{Saeyoung Kim}}, et al., "Development of Multi-Well-Based Electrochemical Dissolved Oxygen Sensor Array," \textit{Sensors and Actuators B: Chemical}, 2020, 306, 127465.
\end{itemize}
\end{sloppypar}

\vfill
\centering{References available upon request.}

\end{document}
